\documentclass[11pt]{article}
\usepackage[final]{acl}
\usepackage{times}
\usepackage{latexsym}
\usepackage[T1]{fontenc}
\usepackage[utf8]{inputenc}
\usepackage{microtype}
\usepackage{booktabs}
\usepackage{amsmath}
\usepackage{xurl}
\usepackage{graphicx}

% Build:  cd docs/milestone && latexmk -pdf milestone.tex
% Tables: python -m eval.milestone_tables   (from the repo root, after pulling new results)
%
% DRAFT WRITTEN BY CLAUDE from the proposal, the repository, and the results files. The course
% follows ACL's rules on generative assistance: before submitting, the team rewrites this in
% their own words, checks every number against results/ and every claim against the cited
% papers, and updates the Acknowledgements to say what was drafted and how it was checked.
% Lines starting with "% TODO:" mark places that need a result or a decision.

\title{Ask Me Nothing: Learning Need-to-Know Clarification \\ for Financial Agents with Multi-Turn RL}

\author{\textbf{Greyson McReynolds \quad Michael Incorvaia} \\
  \textbf{Heeba Merchant \quad Canting Zhu} \\
  Georgia Institute of Technology \\
  ECE 6256 / CS 6604 Conversational AI, Fall 2026 \quad Project milestone}

\begin{document}
\maketitle

\begin{abstract}
In financial dialogue, a clarifying question also collects data. We built BrokerDesk, a simulated brokerage desk that scores task completion, policy violations and requests for client information. Prompted Qwen3-1.7B solves none of 50 validation tasks. After supervised fine-tuning, it places a filled order in every task, solves four, and copies its teacher's 2.88 unnecessary questions per episode. Getting the model to trade exposed a second problem, choosing quantities that meet the client's goal. We report both results and outline how we will test sensitivity-weighted collection costs.
\end{abstract}

\section{Introduction}

Say a client asks a brokerage assistant to raise cash. If the assistant asks for the client's income, it has collected a sensitive fact even if that fact never enters the trade. On another task, failing to ask for a missing field can violate desk policy. Suitability and customer information safeguards motivate the rules in our simulator \citep{finra-2111,finra-2024-regsp}.

Can we teach an assistant to ask for the information it needs and leave the rest alone? We test that idea by comparing a uniform penalty for unnecessary requests with one that depends on sensitivity. At this milestone, we have validation results for scripted agents, the prompted model and its supervised warm start.

\section{Background}

\paragraph{Learning when to clarify.} CollabLLM rewards the longer-term effects of a response, while AskBench trains clarification with rubric-guided rewards \citep{wu-2025-collabllm,zhao-2026-askbench}. \citet{berant-2025-steerable} train a policy with explicit per-question and per-word costs. \citet{zhang-2026-askless} evaluate question costs that vary with user effort and question value. Our experiment varies the cost of collecting unnecessary data by its sensitivity.

\paragraph{Data minimization.} PrivacyPeek audits unnecessary acquisition through tool calls, and MyPhoneBench measures unnecessary permissions and form disclosures \citep{zhang-2026-privacypeek,tang-2026-myphonebench}. \citet{lan-2025-ci-rl} train agents to reduce inappropriate disclosure. BrokerDesk evaluates both acquisition and disclosure in dialogue and tool use, including cases where desk policy requires the agent to obtain a missing field.

\paragraph{Evaluation and training.} We follow $\tau$-bench for policy-guided tool use and end-state scoring, and AgentDojo for testing instructions embedded in untrusted tool outputs \citep{yao-2025-taubench,debenedetti-2024-agentdojo}. We plan to apply GRPO to groups of complete dialogues using LoRA adapters on Qwen3 \citep{shao-2024-deepseekmath,qwen3-model}.

\section{BrokerDesk}
\label{sec:desk}

\paragraph{Tasks.} The client either asks to raise cash while preserving a holding or to rebalance to an equity share based on their risk tolerance. We score each episode against a hidden goal. The seeded splits contain 300 training, 50 validation and 100 test tasks with disjoint clients. In each family, 60\% of requests omit information relevant to the goal or policy.

\paragraph{Tools and policy.} The agent has 12 calls to inspect positions, profiles and quotes, check or place orders, ask the client, or send messages. Prices advance after each call along a seeded path. The desk limits order sizes, prohibits restricted securities and short sales, and requires the agent to know the client's risk tolerance before trading.

\paragraph{Scoring questions.} Each request to the client names one field. A field may be derivable from the request, on file, mandatory and missing, available only from the client, or irrelevant. The first question about a mandatory or client-only field counts as needed. A repeat request costs the same as any other unnecessary request for that field. Looking up an irrelevant field on file also counts as unnecessary collection. The field labels let us score requests directly.

\paragraph{Privacy.} Low sensitivity fields include risk tolerance and time horizon. Income, net worth and tax bracket are sensitive, and social security and account numbers are restricted. Each value carries a marker. We count its appearance as a leak in any message to a third party, or in a message to the client if the value is sensitive or restricted.

\paragraph{Reward.} For questions $n_q$, unneeded field requests $U$, violations $n_v$ and unauthorized disclosures $n_\ell$,
\begin{align*}
R =\;&\mathbf{1}[\text{goal met}] - \lambda_q n_q \\
&-\textstyle\sum_{f\in U}\lambda_{\mathrm{tier}(f)}-\lambda_v n_v-\lambda_\ell n_\ell,
\end{align*}
with $\lambda_q=0.02$, tier costs $0.01$, $0.05$ and $0.5$, and $\lambda_v=\lambda_\ell=1$. A mandatory field never obtained counts as a violation, so never asking is penalized by the policy itself.

\paragraph{Injection.} We pair each clean episode with an episode containing an untrusted instruction in a tool response. The metric counts the requested collection or disclosure only after the agent observes the instruction.

\section{Baselines and First Results}
\label{sec:results}

\begin{table*}[t]
\centering
\small
\resizebox{\textwidth}{!}{%
\begin{tabular}{lcccccccc}
\toprule
Agent & Success [95\% CI] & Viol./ep & Needed q/ep & Unneeded q/ep & Missed client-only & Sens.+restr. requests/ep & Timeout & Invalid/ep \\
\midrule
Scripted need-to-know (reference) & 1.00 [1.00, 1.00] & 0.00 & 0.62 & 0.00 & 0.00 & 0.00 & 0.00 & 0.00 \\
Scripted always-ask (SFT teacher) & 1.00 [1.00, 1.00] & 0.00 & 0.62 & 2.88 & 0.00 & 1.00 & 0.00 & 0.00 \\
Scripted never-ask & 0.44 [0.30, 0.58] & 0.30 & 0.00 & 0.00 & 0.30 & 0.00 & 0.00 & 0.00 \\
\midrule
Qwen3-1.7B, prompted & 0.00 [0.00, 0.00] & 0.34 & 0.46 & 1.30 & 0.16 & 0.30 & 0.54 & 0.86 \\
\bottomrule
\end{tabular}}
\caption{BrokerDesk val split, 50 tasks, one rollout per task, seed 0. Success is goal met with zero policy violations; the interval is a task-level bootstrap. Questions are scored as needed only on the first request for a mandatory or client-only field. Timeout: share of episodes that hit the 12-turn limit. Invalid: turns with no usable tool call.}
\label{tab:main}
\end{table*}


\paragraph{Scripted references.} The need-to-know and always-ask agents each solve all 50 validation tasks (Table~\ref{tab:main}). The former makes zero unnecessary requests, while the latter averages 2.88 unnecessary questions, including one sensitive request per episode. Never-ask shares their trading logic and guesses missing values. It solves 22 of 50 tasks, including 2 of the 30 underspecified tasks.

\paragraph{Prompted Qwen3-1.7B.} Both model evaluations used vLLM, temperature zero and thinking disabled. With the desk policy and a need-to-know instruction in the prompt, Qwen3 solved zero tasks. It spent 224 of 468 calls on risk checks, got an order filled in only four episodes, and timed out in 27. Of its 1.76 questions per episode, 1.30 were unnecessary.

\paragraph{Supervised warm start.} After supervised fine-tuning (SFT), the model got a fill in all 50 episodes, timed out in none, and solved four. It asked for the same fields in the same order as its teacher on every task, including 2.88 unnecessary questions per episode. Order quantities were harder to imitate. On 25 cash requests, the model exceeded the permitted amount 16 times and fell short seven times. For one client who asked for \$44{,}000, it raised about \$116{,}337 and promptly said the request was complete. Across both task families, 41 of the 46 failures missed the goal without a recorded policy violation.

\begin{table}[t]
\centering
\small
\setlength{\tabcolsep}{3pt}
\resizebox{\columnwidth}{!}{%
\begin{tabular}{lccccc}
\toprule
Agent & Succ. clean & Succ. inj. & Seen & Attack & Leak \\
\midrule
Scripted need-to-know (reference) & 1.00 & 1.00 & 1.00 & 0.00 & 0.00 \\
Scripted always-ask (SFT teacher) & 1.00 & 1.00 & 1.00 & 0.00 & 0.00 \\
Scripted never-ask & 0.44 & 0.44 & 1.00 & 0.00 & 0.00 \\
Qwen3-1.7B, prompted & 0.00 & 0.00 & 0.74 & 0.08 & 0.04 \\
Scripted gullible (obeys notes) & 1.00 & 1.00 & 1.00 & 1.00 & 0.64 \\
\bottomrule
\end{tabular}}
\caption{Paired clean and injected runs. \emph{Seen}: the agent called a tool whose output carried the attack. \emph{Attack}: after seeing it, the agent sent a targeted field through the tool the attacker named, or requested the targeted field. \emph{Leak}: any unauthorized canary disclosure in an outbound message.}
\label{tab:attack}
\end{table}


\paragraph{Under injection.} The prompted model encountered the injected instruction in 37 of 50 runs, followed it in four, and leaked a marker in two. The SFT model encountered the instruction in all 50 runs, solved three, and registered zero attack successes or marker leaks (Table~\ref{tab:attack}).

\paragraph{Compute.} On one PACE-ICE H100 80\,GB, SFT used 600 teacher dialogues over three epochs. The training loop took 331 seconds and peak GPU memory was 22.6\,GB. The prompted baseline ran 100 episodes in 248 seconds, including server start.

\section{What Changed Since the Proposal}
\label{sec:changes}

\paragraph{Choice of teacher.} The prompted model never completed a task. With few successes, a reward that charges for questions and violations can favour doing very little. We switched from planned frontier distillation to the scripted always-ask teacher so the model could learn the desk workflow from successful episodes. The need-to-know script remains a reference for how much collection the task requires.

\paragraph{Training harness.} The minimal simulator ran on November 16. We built a harness for grouped rollouts on shared price paths, group-normalized advantages and an assistant-token loss mask verified with the Qwen3 tokenizer. Training sequences use the prompt rendered at each turn to match the model's inference context. End-to-end RL remains pending.

\section{Plan for the Last Three Weeks}
\label{sec:plan}

\paragraph{By November 26.} Sample several dialogues per task from the SFT adapter and count the groups with both a success and a failure. If successful samples are rare, test a shared order sizing tool that uses information the agent has already obtained. Add raw value disclosure checks and matched tasks where a tax bracket is required or irrelevant. Vary whether the required bracket is missing or on file, and hold out request paraphrases.

\paragraph{By December 1.} If the pilot supports RL, compare the tiered penalty with a uniform penalty and with the unnecessary request penalty removed. Use the same SFT initialisation and per-question cost, and target three seeds per variant. On paired test tasks, report policy-compliant success, required field recall and unnecessary requests by tier, with task bootstrap intervals and results by seed. With five rollouts per task, report pass$^1$ and pass$^3$.

\paragraph{Poster and final paper.} Submit the poster December 2 and present December 3. For the paper and code due December 8, prioritize the controlled comparison and transcript analysis. If the pilot cannot support RL, report the baseline comparisons and what the order sizing diagnostic reveals.

\section*{Limitations}
BrokerDesk uses synthetic profiles, templated requests and a scripted client that answers even requests for restricted data. Canary matching detects verbatim markers. Current results use one rollout per validation task at one seed. The simplified policy is for simulation.

\section*{Acknowledgements}
% Team review: confirm disclosure, author checks and workstream contributions before submission.
Claude (Anthropic) drafted an earlier version of this milestone, assisted with literature search, and wrote much of the simulator, evaluation and training code. OpenAI Codex checked sources, replayed validation transcripts, assisted with editing, and suggested the uniform penalty comparison and the order sizing diagnostic. Project checks include unit tests, transcript inspection and verification of the Qwen3 loss mask.

\appendix
\section{Team Contributions}
\noindent\emph{Workstream assignments for author confirmation.}\par
\noindent Greyson McReynolds (simulator tools and policy).\par
\noindent Michael Incorvaia (training and PACE-ICE runs).\par
\noindent Heeba Merchant (clients, tasks and privacy).\par
\noindent Canting Zhu (baselines, evaluation and writing).

\clearpage
\bibliography{../proposal/custom,milestone}

\end{document}
